% Options for packages loaded elsewhere
% Options for packages loaded elsewhere
\PassOptionsToPackage{unicode}{hyperref}
\PassOptionsToPackage{hyphens}{url}
\PassOptionsToPackage{dvipsnames,svgnames,x11names}{xcolor}
%
\documentclass[
  11pt,
]{article}
\usepackage{xcolor}
\usepackage[margin=1in]{geometry}
\usepackage{amsmath,amssymb}
\setcounter{secnumdepth}{5}
\usepackage{iftex}
\ifPDFTeX
  \usepackage[T1]{fontenc}
  \usepackage[utf8]{inputenc}
  \usepackage{textcomp} % provide euro and other symbols
\else % if luatex or xetex
  \usepackage{unicode-math} % this also loads fontspec
  \defaultfontfeatures{Scale=MatchLowercase}
  \defaultfontfeatures[\rmfamily]{Ligatures=TeX,Scale=1}
\fi
\usepackage{lmodern}
\ifPDFTeX\else
  % xetex/luatex font selection
\fi
% Use upquote if available, for straight quotes in verbatim environments
\IfFileExists{upquote.sty}{\usepackage{upquote}}{}
\IfFileExists{microtype.sty}{% use microtype if available
  \usepackage[]{microtype}
  \UseMicrotypeSet[protrusion]{basicmath} % disable protrusion for tt fonts
}{}
\makeatletter
\@ifundefined{KOMAClassName}{% if non-KOMA class
  \IfFileExists{parskip.sty}{%
    \usepackage{parskip}
  }{% else
    \setlength{\parindent}{0pt}
    \setlength{\parskip}{6pt plus 2pt minus 1pt}}
}{% if KOMA class
  \KOMAoptions{parskip=half}}
\makeatother
% Make \paragraph and \subparagraph free-standing
\makeatletter
\ifx\paragraph\undefined\else
  \let\oldparagraph\paragraph
  \renewcommand{\paragraph}{
    \@ifstar
      \xxxParagraphStar
      \xxxParagraphNoStar
  }
  \newcommand{\xxxParagraphStar}[1]{\oldparagraph*{#1}\mbox{}}
  \newcommand{\xxxParagraphNoStar}[1]{\oldparagraph{#1}\mbox{}}
\fi
\ifx\subparagraph\undefined\else
  \let\oldsubparagraph\subparagraph
  \renewcommand{\subparagraph}{
    \@ifstar
      \xxxSubParagraphStar
      \xxxSubParagraphNoStar
  }
  \newcommand{\xxxSubParagraphStar}[1]{\oldsubparagraph*{#1}\mbox{}}
  \newcommand{\xxxSubParagraphNoStar}[1]{\oldsubparagraph{#1}\mbox{}}
\fi
\makeatother

\usepackage{color}
\usepackage{fancyvrb}
\newcommand{\VerbBar}{|}
\newcommand{\VERB}{\Verb[commandchars=\\\{\}]}
\DefineVerbatimEnvironment{Highlighting}{Verbatim}{commandchars=\\\{\}}
% Add ',fontsize=\small' for more characters per line
\usepackage{framed}
\definecolor{shadecolor}{RGB}{241,243,245}
\newenvironment{Shaded}{\begin{snugshade}}{\end{snugshade}}
\newcommand{\AlertTok}[1]{\textcolor[rgb]{0.68,0.00,0.00}{#1}}
\newcommand{\AnnotationTok}[1]{\textcolor[rgb]{0.37,0.37,0.37}{#1}}
\newcommand{\AttributeTok}[1]{\textcolor[rgb]{0.40,0.45,0.13}{#1}}
\newcommand{\BaseNTok}[1]{\textcolor[rgb]{0.68,0.00,0.00}{#1}}
\newcommand{\BuiltInTok}[1]{\textcolor[rgb]{0.00,0.23,0.31}{#1}}
\newcommand{\CharTok}[1]{\textcolor[rgb]{0.13,0.47,0.30}{#1}}
\newcommand{\CommentTok}[1]{\textcolor[rgb]{0.37,0.37,0.37}{#1}}
\newcommand{\CommentVarTok}[1]{\textcolor[rgb]{0.37,0.37,0.37}{\textit{#1}}}
\newcommand{\ConstantTok}[1]{\textcolor[rgb]{0.56,0.35,0.01}{#1}}
\newcommand{\ControlFlowTok}[1]{\textcolor[rgb]{0.00,0.23,0.31}{\textbf{#1}}}
\newcommand{\DataTypeTok}[1]{\textcolor[rgb]{0.68,0.00,0.00}{#1}}
\newcommand{\DecValTok}[1]{\textcolor[rgb]{0.68,0.00,0.00}{#1}}
\newcommand{\DocumentationTok}[1]{\textcolor[rgb]{0.37,0.37,0.37}{\textit{#1}}}
\newcommand{\ErrorTok}[1]{\textcolor[rgb]{0.68,0.00,0.00}{#1}}
\newcommand{\ExtensionTok}[1]{\textcolor[rgb]{0.00,0.23,0.31}{#1}}
\newcommand{\FloatTok}[1]{\textcolor[rgb]{0.68,0.00,0.00}{#1}}
\newcommand{\FunctionTok}[1]{\textcolor[rgb]{0.28,0.35,0.67}{#1}}
\newcommand{\ImportTok}[1]{\textcolor[rgb]{0.00,0.46,0.62}{#1}}
\newcommand{\InformationTok}[1]{\textcolor[rgb]{0.37,0.37,0.37}{#1}}
\newcommand{\KeywordTok}[1]{\textcolor[rgb]{0.00,0.23,0.31}{\textbf{#1}}}
\newcommand{\NormalTok}[1]{\textcolor[rgb]{0.00,0.23,0.31}{#1}}
\newcommand{\OperatorTok}[1]{\textcolor[rgb]{0.37,0.37,0.37}{#1}}
\newcommand{\OtherTok}[1]{\textcolor[rgb]{0.00,0.23,0.31}{#1}}
\newcommand{\PreprocessorTok}[1]{\textcolor[rgb]{0.68,0.00,0.00}{#1}}
\newcommand{\RegionMarkerTok}[1]{\textcolor[rgb]{0.00,0.23,0.31}{#1}}
\newcommand{\SpecialCharTok}[1]{\textcolor[rgb]{0.37,0.37,0.37}{#1}}
\newcommand{\SpecialStringTok}[1]{\textcolor[rgb]{0.13,0.47,0.30}{#1}}
\newcommand{\StringTok}[1]{\textcolor[rgb]{0.13,0.47,0.30}{#1}}
\newcommand{\VariableTok}[1]{\textcolor[rgb]{0.07,0.07,0.07}{#1}}
\newcommand{\VerbatimStringTok}[1]{\textcolor[rgb]{0.13,0.47,0.30}{#1}}
\newcommand{\WarningTok}[1]{\textcolor[rgb]{0.37,0.37,0.37}{\textit{#1}}}

\usepackage{longtable,booktabs,array}
\newcounter{none} % for unnumbered tables
\usepackage{calc} % for calculating minipage widths
% Correct order of tables after \paragraph or \subparagraph
\usepackage{etoolbox}
\makeatletter
\patchcmd\longtable{\par}{\if@noskipsec\mbox{}\fi\par}{}{}
\makeatother
% Allow footnotes in longtable head/foot
\IfFileExists{footnotehyper.sty}{\usepackage{footnotehyper}}{\usepackage{footnote}}
\makesavenoteenv{longtable}
\usepackage{graphicx}
\makeatletter
\newsavebox\pandoc@box
\newcommand*\pandocbounded[1]{% scales image to fit in text height/width
  \sbox\pandoc@box{#1}%
  \Gscale@div\@tempa{\textheight}{\dimexpr\ht\pandoc@box+\dp\pandoc@box\relax}%
  \Gscale@div\@tempb{\linewidth}{\wd\pandoc@box}%
  \ifdim\@tempb\p@<\@tempa\p@\let\@tempa\@tempb\fi% select the smaller of both
  \ifdim\@tempa\p@<\p@\scalebox{\@tempa}{\usebox\pandoc@box}%
  \else\usebox{\pandoc@box}%
  \fi%
}
% Set default figure placement to htbp
\def\fps@figure{htbp}
\makeatother





\setlength{\emergencystretch}{3em} % prevent overfull lines

\providecommand{\tightlist}{%
  \setlength{\itemsep}{0pt}\setlength{\parskip}{0pt}}



 


\usepackage{amsmath,amssymb,amsthm,bm}
\usepackage{mathtools}
\makeatletter
\@ifpackageloaded{float}{}{\usepackage{float}}
\floatstyle{plain}
\@ifundefined{c@chapter}{\newfloat{assumption}{h}{loasm}}{\newfloat{assumption}{h}{loasm}[chapter]}
\floatname{assumption}{Assumption}
\newcommand*\listofassumptions{\listof{assumption}{List of Assumptions}}
\makeatother
\makeatletter
\@ifpackageloaded{float}{}{\usepackage{float}}
\floatstyle{plain}
\@ifundefined{c@chapter}{\newfloat{remark}{h}{lorem}}{\newfloat{remark}{h}{lorem}[chapter]}
\floatname{remark}{Remark}
\newcommand*\listofremarks{\listof{remark}{List of Remarks}}
\makeatother
\makeatletter
\@ifpackageloaded{tcolorbox}{}{\usepackage[skins,breakable]{tcolorbox}}
\@ifpackageloaded{fontawesome5}{}{\usepackage{fontawesome5}}
\definecolor{quarto-callout-color}{HTML}{909090}
\definecolor{quarto-callout-note-color}{HTML}{0758E5}
\definecolor{quarto-callout-important-color}{HTML}{CC1914}
\definecolor{quarto-callout-warning-color}{HTML}{EB9113}
\definecolor{quarto-callout-tip-color}{HTML}{00A047}
\definecolor{quarto-callout-caution-color}{HTML}{FC5300}
\definecolor{quarto-callout-color-frame}{HTML}{acacac}
\definecolor{quarto-callout-note-color-frame}{HTML}{4582ec}
\definecolor{quarto-callout-important-color-frame}{HTML}{d9534f}
\definecolor{quarto-callout-warning-color-frame}{HTML}{f0ad4e}
\definecolor{quarto-callout-tip-color-frame}{HTML}{02b875}
\definecolor{quarto-callout-caution-color-frame}{HTML}{fd7e14}
\makeatother
\makeatletter
\@ifpackageloaded{caption}{}{\usepackage{caption}}
\AtBeginDocument{%
\ifdefined\contentsname
  \renewcommand*\contentsname{Table of contents}
\else
  \newcommand\contentsname{Table of contents}
\fi
\ifdefined\listfigurename
  \renewcommand*\listfigurename{List of Figures}
\else
  \newcommand\listfigurename{List of Figures}
\fi
\ifdefined\listtablename
  \renewcommand*\listtablename{List of Tables}
\else
  \newcommand\listtablename{List of Tables}
\fi
\ifdefined\figurename
  \renewcommand*\figurename{Fig.}
\else
  \newcommand\figurename{Fig.}
\fi
\ifdefined\tablename
  \renewcommand*\tablename{Table}
\else
  \newcommand\tablename{Table}
\fi
}
\@ifpackageloaded{float}{}{\usepackage{float}}
\floatstyle{ruled}
\@ifundefined{c@chapter}{\newfloat{codelisting}{h}{lop}}{\newfloat{codelisting}{h}{lop}[chapter]}
\floatname{codelisting}{Listing}
\newcommand*\listoflistings{\listof{codelisting}{List of Listings}}
\usepackage{amsthm}
\theoremstyle{plain}
\newtheorem{lemma}{Lemma}[section]
\theoremstyle{plain}
\newtheorem{theorem}{Theorem}[section]
\theoremstyle{definition}
\newtheorem{definition}{Definition}[section]
\theoremstyle{remark}
\AtBeginDocument{\renewcommand*{\proofname}{Proof}}
\newtheorem*{remark}{Remark}
\newtheorem*{solution}{Solution}
\newtheorem{refremark}{Remark}[section]
\newtheorem{refsolution}{Solution}[section]
\makeatother
\makeatletter
\makeatother
\makeatletter
\@ifpackageloaded{caption}{}{\usepackage{caption}}
\@ifpackageloaded{subcaption}{}{\usepackage{subcaption}}
\makeatother
\usepackage{bookmark}
\IfFileExists{xurl.sty}{\usepackage{xurl}}{} % add URL line breaks if available
\urlstyle{same}
\hypersetup{
  pdftitle={Stability-Aware Residual Reinforcement Learning},
  pdfauthor={Piper Manipulator Project},
  colorlinks=true,
  linkcolor={blue},
  filecolor={Maroon},
  citecolor={Blue},
  urlcolor={Blue},
  pdfcreator={LaTeX via pandoc}}


\title{Stability-Aware Residual Reinforcement Learning}
\usepackage{etoolbox}
\makeatletter
\providecommand{\subtitle}[1]{% add subtitle to \maketitle
  \apptocmd{\@title}{\par {\large #1 \par}}{}{}
}
\makeatother
\subtitle{Input-to-State Stability Analysis and Action-Clipping Law\\
for a Cascade NMPC--DOB--RL Controller on a 6-DOF Manipulator}
\author{Piper Manipulator Project}
\date{2026-05-29}
\begin{document}
\maketitle

\renewcommand*\contentsname{Table of contents}
{
\setcounter{tocdepth}{3}
\tableofcontents
}

\section{Introduction}\label{sec-intro}

This report establishes a deterministic stability guarantee for a
residual reinforcement-learning (RL) controller deployed on top of a
nominal nonlinear model-predictive controller (NMPC) and a model-based
disturbance observer (DOB). The combined NMPC--DOB layer is mature and
admits a classical stability certificate; the RL layer, by contrast, is
a neural-network policy and resists direct stability analysis.

Our strategy is to \emph{borrow} the stability certificate of the
baseline. We first show that the NMPC--DOB cascade is
\emph{Input-to-State Stable} (ISS) with respect to the uncompensated
residual disturbance that escapes the observer's passband. From the ISS
inequality we extract a quantitative budget on the admissible residual,
then \emph{map this budget back onto the RL action space} as a hard,
state-aware magnitude ceiling. Any RL torque within the ceiling is
admissible by construction; any torque exceeding the ceiling is radially
scaled back. The cage is built entirely from quantities of the certified
baseline, so the learning agent inherits a provable stability margin
without itself being verified.

The remainder of the document is organised as follows.
Section~\ref{sec-setup} fixes notation and the cascade control law.
Section~\ref{sec-cascade} introduces the three layers.
Section~\ref{sec-lyap} states the standing assumptions that make the
NMPC optimal cost a valid Lyapunov candidate. Section~\ref{sec-iss}
derives the baseline ISS inequality. Section~\ref{sec-clip} is the core
contribution: the ISS-based action-clipping law and its derivation via
the Cauchy--Schwarz inequality. Section~\ref{sec-estim} describes how
the three constants \(c_x\), \(\gamma_0\), and \(r\) are obtained from
simulation and hardware experiments.

\section{Setup and Notation}\label{sec-setup}

\subsection{Tracking-error state and the cascade control
law}\label{tracking-error-state-and-the-cascade-control-law}

We work in the discrete-time tracking-error coordinates of the 6-DOF
manipulator. Let \(\bm{q}_k,\bm{v}_k \in \mathbb{R}^6\) denote the
measured joint position and velocity at sample \(k\), and let
\(\bm{q}^{\mathrm{ref}}_k,\bm{v}^{\mathrm{ref}}_k\) be the time-varying
joint-space reference obtained by inverse kinematics of the Cartesian
end-effector target. The tracking-error state is

\begin{equation}\protect\phantomsection\label{eq-state}{
\bm{x}_k \;\triangleq\;
\begin{bmatrix}
\bm{q}_k - \bm{q}^{\mathrm{ref}}_k \\[2pt]
\bm{v}_k - \bm{v}^{\mathrm{ref}}_k
\end{bmatrix}
\;\in\; \mathbb{R}^{12}.
}\end{equation}

The origin \(\bm{x}_k = \bm{0}\) denotes perfect tracking. Throughout
this document, \(\|\,\cdot\,\|\) denotes the Euclidean norm.

The cascade control law executed by our implementation is

\begin{equation}\protect\phantomsection\label{eq-cmd}{
\bm{u}_{\mathrm{cmd},k}
\;=\;
\underbrace{\bm{u}_{\mathrm{MPC},k}}_{\text{nominal NMPC}}
\;-\;
\underbrace{\bm{d}_{\mathrm{filter},k}}_{\text{DOB estimate}}
\;-\;
\underbrace{\bm{u}_{\mathrm{RL},k}}_{\text{residual RL}},
}\end{equation}

where all three components live in the same physical joint-torque space
\(\mathbb{R}^6\). The minus sign on the RL term reflects that the policy
is trained to \emph{cancel} the residual disturbance, mirroring the DOB.

\subsection{Plant dynamics under the residual
disturbance}\label{plant-dynamics-under-the-residual-disturbance}

Let \(\bm{f}(\bm{x},\bm{u})\) denote the nominal (URDF-modelled)
closed-loop dynamics. The true plant, on real hardware or in simulation
with friction and base excitation, evolves as

\begin{equation}\protect\phantomsection\label{eq-plant}{
\bm{x}_{k+1} \;=\; \bm{f}\bigl(\bm{x}_k,\bm{u}_{\mathrm{cmd},k}\bigr)
\;+\; \bm{w}_k,
}\end{equation}

where \(\bm{w}_k\) collects all sources of model mismatch (un-modelled
friction, base excitation, payload variation, parameter drift). The DOB
estimates the in-band component of \(\bm{w}_k\) and subtracts it as
\(\bm{d}_{\mathrm{filter},k}\); what the observer cannot capture remains
as the \emph{residual disturbance}

\begin{equation}\protect\phantomsection\label{eq-dres-def}{
\bm{d}_{\mathrm{res},k} \;\triangleq\;
\bm{w}_k \;-\; \bm{d}_{\mathrm{filter},k}
\;\in\;\mathbb{R}^6.
}\end{equation}

The residual is the \emph{only} exogenous input that appears in the
forthcoming ISS analysis. We will frequently abuse notation and write
\(\bm{x}_{k+1} = \bm{x}_{k+1}^{\mathrm{nom}} + \bm{d}_{\mathrm{res},k}\),
where \(\bm{x}_{k+1}^{\mathrm{nom}}\) is the disturbance-free successor
state under the nominal closed loop; this is the standard
state-displacement embedding of an input disturbance and is sufficient
for the Lyapunov analysis that follows.

\section{The Cascade Architecture}\label{sec-cascade}

\subsection{Layer 1 --- Nominal NMPC}\label{sec-nmpc}

The NMPC solves, at each sample \(k\), a finite-horizon optimal-control
problem

\begin{equation}\protect\phantomsection\label{eq-vstar}{
V_N^*(\bm{x}_k)
\;=\;
\min_{\bm{u}_{0:N-1}}
\sum_{j=0}^{N-1} \ell\bigl(\bm{x}_{k+j},\bm{u}_{k+j}\bigr)
\;+\; V_f\bigl(\bm{x}_{k+N}\bigr)
}\end{equation}

subject to the nominal dynamics
\(\bm{x}_{j+1}=\bm{f}(\bm{x}_j,\bm{u}_j)\), input bounds, and
(implicitly) a terminal set. The NMPC feedback law is the first element
of the minimising sequence,
\(\bm{u}_{\mathrm{MPC},k} = \kappa_N(\bm{x}_k) \triangleq
\bm{u}^*_{0}(\bm{x}_k)\). In our implementation the stage cost is
\begin{equation}\protect\phantomsection\label{eq-stage}{
\ell(\bm{x},\bm{u}) \;=\;
Q_{\mathrm{pos}}\,\bigl\|\bm{p}(\bm{x})-\bm{p}^{\star}\bigr\|^2
+ Q_{\mathrm{vel}}\,\bigl\|\bm{v}(\bm{x})\bigr\|^2
+ R_{\tau}\,\|\bm{u}\|^2,
}\end{equation} with \(Q_{\mathrm{pos}}=5000\), \(Q_{\mathrm{vel}}=10\),
\(R_{\tau}=0.1\), horizon \(N=8\), and sample time
\(\Delta t = 0.02\)\textasciitilde s; \(\bm{p}(\bm{x})\) denotes the
forward-kinematics end-effector position.

The optimal cost \(V_N^*\) defined in Eq.~\ref{eq-vstar} will serve as
our Lyapunov candidate.

\subsection{Layer 2 --- Model-based DOB}\label{sec-dob}

The DOB estimates the disturbance from the dynamics residual:
\(\bm{d}_{\mathrm{raw},k} = \mathrm{RNEA}\bigl(\bm{q}_k,\bm{v}_k,
\hat{\bm{a}}_k\bigr) - \bm{u}_{\mathrm{cmd},k-1}\), where
\(\hat{\bm{a}}_k\) is an acceleration estimate (in simulation this is
the ground-truth physical acceleration; on hardware it is the finite
difference of \(\bm{v}_k\)). The raw estimate is passed through a
first-order low-pass filter with smoothing factor \(\alpha \in (0,1)\):

\begin{equation}\protect\phantomsection\label{eq-dob-lpf}{
\bm{d}_{\mathrm{filter},k}
\;=\; (1-\alpha)\,\bm{d}_{\mathrm{filter},k-1}
\;+\; \alpha\,\bm{d}_{\mathrm{raw},k}.
}\end{equation}

The filter shapes which disturbances are cancelled and which fall
through to the residual Eq.~\ref{eq-dres-def}. Reducing \(\alpha\)
tightens the passband, which \emph{enlarges} the residual that the RL
layer must handle.

\subsection{Layer 3 --- Residual RL}\label{sec-rl}

The RL policy \(\bm{u}_{\mathrm{RL},k} = \pi(\bm{o}_k)\) is a neural
network that observes the same signals available to the DOB plus the
DOB's output, and emits an additive joint torque whose ideal target is
the residual itself, i.e.
\(\bm{u}_{\mathrm{RL},k} \to \bm{d}_{\mathrm{res},k}\). The policy is
trained against a reward that penalises the residual
\(\bm{d}_{\mathrm{eff},k} \triangleq \bm{d}_{\mathrm{res},k} -
\bm{u}_{\mathrm{RL},k}\) and tracking error. Because \(\pi\) is a
learned, possibly stochastic, function, no closed-form stability
statement is available for the closed loop that includes it. The role of
Section~\ref{sec-clip} is to provide one.

\section{Lyapunov Candidate: The Optimal NMPC Cost}\label{sec-lyap}

We adopt the NMPC value function \(V_N^*\) from Eq.~\ref{eq-vstar} as
the composite Lyapunov candidate for the cascade. This choice is
standard in the predictive-control literature and avoids the need to
design a Lyapunov function by hand.

\subsection{Standing assumptions}\label{standing-assumptions}

We collect the three structural assumptions used throughout the rest of
the document. They are the standard hypotheses behind the quasi-infinite
horizon NMPC stability theorem {[}@mayne2000; @rawlings2017{]}, adapted
to our cascade setting.

\begin{assumption}

\centering{

\textbf{Standing assumptions on the operating domain and the NMPC value
function.}

There exists a compact operating domain \(\mathcal{X} \subset
\mathbb{R}^{12}\) containing the origin in its interior, such that on
\(\mathcal{X}\) the following hold:

\begin{itemize}
\item
  \textbf{(A1) Lyapunov sandwich.} There exist \(\alpha_1,\alpha_2 \in
  \mathcal{K}_\infty\) such that
  \(\alpha_1(\|\bm{x}\|) \le V_N^*(\bm{x}) \le \alpha_2(\|\bm{x}\|)\)
  for all \(\bm{x}\in\mathcal{X}\).
\item
  \textbf{(A2) Quadratic stage-cost lower bound.} There exists \(c_x>0\)
  with
  \(\ell\bigl(\bm{x},\kappa_N(\bm{x})\bigr) \;\ge\; c_x\,\|\bm{x}\|^2\)
  for all \(\bm{x}\in\mathcal{X}\).
\item
  \textbf{(A3) Local Lipschitz continuity of \(V_N^*\).} There exists
  \(\gamma_0>0\) such that, for every \(\bm{x}\in\mathcal{X}\) and every
  perturbation \(\bm{\delta}\in\mathbb{R}^{12}\) with
  \(\bm{x}+\bm{\delta}\in\mathcal{X}\),
  \(\bigl|V_N^*(\bm{x}+\bm{\delta}) - V_N^*(\bm{x})\bigr|
  \;\le\; \gamma_0\,\|\bm{\delta}\|\).
\end{itemize}

}

\caption{\label{asm-lyap}}

\end{assumption}%

(A1)--(A3) are mild: (A1) follows from positive-definiteness and radial
unboundedness of the cost; (A2) is design-driven and is automatically
satisfied by the quadratic stage cost Eq.~\ref{eq-stage}; (A3) holds
locally because \(V_N^*\) is the value of a smooth, finite-horizon
quadratic program on a compact set.

The constants \(c_x\) and \(\gamma_0\) that appear in (A2)--(A3) are
\emph{exactly} the two constants we will need to size the
action-clipping law. How to obtain them in practice is the subject of
Section~\ref{sec-estim}.

\subsection{Value-function descent (nominal
stability)}\label{sec-vfdescent}

Under (A1)--(A2) and the standard NMPC terminal ingredients (a
control-invariant terminal set \(\mathcal{X}_f \subset \mathcal{X}\)
with a local control-Lyapunov terminal cost \(V_f\)), the NMPC value
function satisfies the well-known \emph{descent property} along the
nominal closed loop.

\begin{lemma}[]\protect\hypertarget{lem-descent}{}\label{lem-descent}

\textbf{Lemma 1 (Value-function descent).} Under
Assumption~\ref{asm-lyap} and the standard NMPC terminal assumptions,
the optimal cost \(V_N^*\) satisfies
\begin{equation}\protect\phantomsection\label{eq-descent}{
V_N^*\bigl(\bm{x}_{k+1}^{\mathrm{nom}}\bigr) - V_N^*(\bm{x}_k)
\;\le\;
-\ell\bigl(\bm{x}_k,\kappa_N(\bm{x}_k)\bigr)
\;\le\;
-c_x\,\|\bm{x}_k\|^2
}\end{equation} for every \(\bm{x}_k \in \mathcal{X}\), where
\(\bm{x}_{k+1}^{\mathrm{nom}} = \bm{f}\bigl(\bm{x}_k,\kappa_N(\bm{x}_k)\bigr)\)
is the disturbance-free successor.

\end{lemma}

\begin{tcolorbox}[enhanced jigsaw, arc=.35mm, bottomrule=.15mm, bottomtitle=1mm, breakable, colback=white, colbacktitle=quarto-callout-note-color!10!white, colframe=quarto-callout-note-color-frame, coltitle=black, left=2mm, leftrule=.75mm, opacityback=0, opacitybacktitle=0.6, rightrule=.15mm, title=\textcolor{quarto-callout-note-color}{\faInfo}\hspace{0.5em}{Proof sketch (standard)}, titlerule=0mm, toprule=.15mm, toptitle=1mm]

The argument is the canonical ``tail-shift'' of NMPC stability theory
{[}@mayne2000, Thm. 1; @rawlings2017, §2.4{]}. From \(\bm{x}_k\) the
optimal plan costs \(V_N^*(\bm{x}_k)\) over \(N\) steps. After applying
its first input the system reaches \(\bm{x}_{k+1}^{\mathrm{nom}}\);
\emph{re-using the tail} of the previous plan and appending the terminal
controller yields a \emph{feasible} (not necessarily optimal) plan from
\(\bm{x}_{k+1}^{\mathrm{nom}}\) whose cost is at most
\(V_N^*(\bm{x}_k) - \ell(\bm{x}_k,\kappa_N(\bm{x}_k))\). The
re-optimised value \(V_N^*(\bm{x}_{k+1}^{\mathrm{nom}})\) cannot exceed
this feasible cost, which proves the first inequality. Applying (A2)
proves the second.

\end{tcolorbox}

Lemma~\ref{lem-descent} isolates the \emph{intrinsic convergence
strength} of the controller: every disturbance-free step strips at least
\(c_x\|\bm{x}_k\|^2\) units of cost. The constant \(c_x\) is the
\textbf{convergence coefficient} of the nominal loop and is the first
ingredient of the action-clipping law.

\section{Input-to-State Stability of the Baseline}\label{sec-iss}

We now activate the residual disturbance and the DOB, keeping
\(\bm{u}_{\mathrm{RL},k} = \bm{0}\) (i.e., the \emph{baseline} is the
NMPC--DOB cascade alone). The goal is to certify ISS of the baseline
with respect to \(\bm{d}_{\mathrm{res},k}\).

\subsection{Cost propagation under the residual
disturbance}\label{cost-propagation-under-the-residual-disturbance}

The true successor state under the residual disturbance is
\begin{equation}\protect\phantomsection\label{eq-disturbed-succ}{
\bm{x}_{k+1} \;=\; \bm{x}_{k+1}^{\mathrm{nom}}
\;+\; \bm{d}_{\mathrm{res},k}.
}\end{equation}

The one-step change of the Lyapunov candidate decomposes by adding and
subtracting \(V_N^*(\bm{x}_{k+1}^{\mathrm{nom}})\):

\begin{equation}\protect\phantomsection\label{eq-split}{
\begin{aligned}
V_N^*(\bm{x}_{k+1}) - V_N^*(\bm{x}_k)
&= \underbrace{\bigl[V_N^*(\bm{x}_{k+1}^{\mathrm{nom}}
                 + \bm{d}_{\mathrm{res},k})
        - V_N^*(\bm{x}_{k+1}^{\mathrm{nom}})\bigr]}_{\text{perturbation term}}
\\[4pt]
&\quad+ \underbrace{\bigl[V_N^*(\bm{x}_{k+1}^{\mathrm{nom}})
        - V_N^*(\bm{x}_k)\bigr]}_{\text{nominal descent term}}.
\end{aligned}
}\end{equation}

This decomposition is an identity: nothing has been bounded yet. Its
purpose is to \emph{isolate the only place the disturbance enters}---the
first bracket---from the part already controlled by
Lemma~\ref{lem-descent}.

\subsection{Lipschitz bound on the perturbation
term}\label{lipschitz-bound-on-the-perturbation-term}

The first bracket of Eq.~\ref{eq-split} is bounded directly by
Assumption (A3): \begin{equation}\protect\phantomsection\label{eq-lip}{
V_N^*\bigl(\bm{x}_{k+1}^{\mathrm{nom}} + \bm{d}_{\mathrm{res},k}\bigr)
- V_N^*\bigl(\bm{x}_{k+1}^{\mathrm{nom}}\bigr)
\;\le\; \gamma_0\,\bigl\|\bm{d}_{\mathrm{res},k}\bigr\|.
}\end{equation}

The constant \(\gamma_0\) is the \textbf{disturbance gain} of the
cascade. It plays the role of an ``exchange rate'' from units of state
displacement (metres of error) to units of cost. A large \(\gamma_0\)
means the cost landscape is steep near the operating trajectory and the
system is sensitive to small disturbances; a small \(\gamma_0\) means
the cost surface is flat and the system is tolerant.

\subsection{The baseline ISS
inequality}\label{the-baseline-iss-inequality}

Combining Lemma~\ref{lem-descent} and Eq.~\ref{eq-lip} yields the
central result of this section.

\begin{theorem}[]\protect\hypertarget{thm-baseline-iss}{}\label{thm-baseline-iss}

\textbf{Theorem 1 (Baseline ISS).} Under Assumption~\ref{asm-lyap} and
the standard NMPC terminal assumptions, the NMPC--DOB baseline cascade
(i.e., Eq.~\ref{eq-cmd} with \(\bm{u}_{\mathrm{RL}}=\bm{0}\)) satisfies,
on the operating domain \(\mathcal{X}\),
\begin{equation}\protect\phantomsection\label{eq-iss}{
V_N^*(\bm{x}_{k+1}) - V_N^*(\bm{x}_k)
\;\le\;
-c_x\,\|\bm{x}_k\|^2
\;+\; \gamma_0\,\|\bm{d}_{\mathrm{res},k}\|.
}\end{equation} Equivalently, the baseline is \emph{Input-to-State
Stable} with respect to the residual disturbance
\(\bm{d}_{\mathrm{res},k}\), with ISS-Lyapunov function \(V_N^*\),
dissipation rate \(\alpha_3(s)=c_x s^2\), and disturbance gain
\(\sigma(s)=\gamma_0 s\).

\end{theorem}

\begin{tcolorbox}[enhanced jigsaw, arc=.35mm, bottomrule=.15mm, bottomtitle=1mm, breakable, colback=white, colbacktitle=quarto-callout-note-color!10!white, colframe=quarto-callout-note-color-frame, coltitle=black, left=2mm, leftrule=.75mm, opacityback=0, opacitybacktitle=0.6, rightrule=.15mm, title=\textcolor{quarto-callout-note-color}{\faInfo}\hspace{0.5em}{Proof}, titlerule=0mm, toprule=.15mm, toptitle=1mm]

The decomposition Eq.~\ref{eq-split} is an identity. Bound its first
bracket using Eq.~\ref{eq-lip} and its second using
Eq.~\ref{eq-descent}; sum the two bounds. The resulting inequality
Eq.~\ref{eq-iss} matches the ISS-Lyapunov template of Jiang \& Wang
{[}@jiang2001, Thm. 1{]}, so ISS follows immediately. \(\square\)

\end{tcolorbox}

Eq.~\ref{eq-iss} is the \emph{quantitative} statement of stability. The
first term on the right is the energy \emph{drain} due to the
controller; it scales like \(\|\bm{x}_k\|^2\) and dominates far from the
origin. The second term is the energy \emph{refill} due to the
disturbance; it scales linearly with \(\|\bm{d}_{\mathrm{res},k}\|\).
Equilibrium between drain and refill sets the steady-state radius of the
tracking error.

\section{ISS-Based Action Clipping}\label{sec-clip}

We now use the budget hidden inside Eq.~\ref{eq-iss} to constrain the RL
action. The argument proceeds in four steps: (i) extract the admissible
residual budget from the ISS inequality; (ii) write the ISS inequality
for the closed loop \emph{with} the RL layer; (iii) impose a
\emph{preservation criterion} that the RL action must not enlarge the
residual it is meant to shrink; (iv) reduce that criterion, via the
Cauchy--Schwarz inequality, to a hard magnitude ceiling on
\(\bm{u}_{\mathrm{RL},k}\).

\subsection{The admissible residual budget}\label{sec-budget}

Fix a desired \emph{practical convergence radius} \(r > 0\): we want the
tracking error to settle inside \(\|\bm{x}_k\| \le r\). From
Eq.~\ref{eq-iss}, the cost stops contracting precisely when the drain no
longer dominates the refill, i.e.~when
\(c_x\,\|\bm{x}_k\|^2 \le \gamma_0\,\|\bm{d}_{\mathrm{res},k}\|\). To
force this balance to occur \emph{at} or \emph{inside} the target
radius, the refill at \(\|\bm{x}_k\|=r\) must not exceed the drain
there: \begin{equation}\protect\phantomsection\label{eq-budget}{
\gamma_0\,\|\bm{d}_{\mathrm{res},k}\| \;\le\; c_x\,r^2
\quad\Longleftrightarrow\quad
\|\bm{d}_{\mathrm{res},k}\| \;\le\; \bar{d}_{\mathrm{res}}
\;\triangleq\; \frac{c_x\,r^2}{\gamma_0}.
}\end{equation}

\begin{definition}[]\protect\hypertarget{def-budget}{}\label{def-budget}

\textbf{Definition 1 (Admissible residual budget).} The scalar
\(\bar{d}_{\mathrm{res}} \triangleq c_x r^2 / \gamma_0\) is the
\textbf{admissible residual budget}: the largest residual norm that the
baseline NMPC--DOB cascade can absorb while keeping the tracking-error
state within radius \(r\).

\end{definition}

The three quantities entering Eq.~\ref{eq-budget} have crisp roles:
\(c_x\) rewards strong convergence (larger \(c_x\) enlarges the budget),
\(r\) rewards looser performance targets (the budget grows like
\(r^2\)), and \(\gamma_0\) punishes sensitivity to disturbance (a
steeper cost surface shrinks the budget).

\subsection{ISS inequality with the RL layer
active}\label{iss-inequality-with-the-rl-layer-active}

With \(\bm{u}_{\mathrm{RL},k} \ne \bm{0}\), the effective disturbance
seen by the baseline is shifted:
\begin{equation}\protect\phantomsection\label{eq-deff}{
\bm{d}_{\mathrm{eff},k}
\;\triangleq\;
\bm{d}_{\mathrm{res},k} - \bm{u}_{\mathrm{RL},k},
}\end{equation} and the ISS inequality Eq.~\ref{eq-iss} becomes
\begin{equation}\protect\phantomsection\label{eq-iss-with-rl}{
V_N^*(\bm{x}_{k+1}) - V_N^*(\bm{x}_k)
\;\le\;
-c_x\,\|\bm{x}_k\|^2
\;+\; \gamma_0\,\|\bm{d}_{\mathrm{res},k} - \bm{u}_{\mathrm{RL},k}\|.
}\end{equation}

\subsection{The stability-preservation
criterion}\label{the-stability-preservation-criterion}

We declare that the RL layer is \emph{admissible at step \(k\)} whenever
it does not enlarge the residual it is intended to cancel:
\begin{equation}\protect\phantomsection\label{eq-criterion}{
\boxed{\quad
\|\bm{d}_{\mathrm{res},k} - \bm{u}_{\mathrm{RL},k}\|
\;\le\; \|\bm{d}_{\mathrm{res},k}\|.
\quad}
}\end{equation}

If Eq.~\ref{eq-criterion} holds, then Eq.~\ref{eq-iss-with-rl} is
\emph{no worse} than Eq.~\ref{eq-iss}, so the baseline ISS certificate
carries over verbatim and the closed loop remains ISS with the same
constants.

\subsection{From the preservation criterion to a magnitude
bound}\label{sec-magbound}

Eq.~\ref{eq-criterion} is geometric and depends on the unmeasurable
internal quantity \(\bm{d}_{\mathrm{res},k}\). We now reduce it to a
\emph{measurable} condition on \(\|\bm{u}_{\mathrm{RL},k}\|\) alone, in
three algebraic steps.

\textbf{Step 1 --- Square both sides.} Squaring Eq.~\ref{eq-criterion}
and expanding the inner product on the left: \[
\|\bm{d}_{\mathrm{res},k}\|^2
- 2\,\bm{d}_{\mathrm{res},k}^{\!\top}\bm{u}_{\mathrm{RL},k}
+ \|\bm{u}_{\mathrm{RL},k}\|^2
\;\le\;
\|\bm{d}_{\mathrm{res},k}\|^2,
\] which simplifies to
\begin{equation}\protect\phantomsection\label{eq-squared}{
\|\bm{u}_{\mathrm{RL},k}\|^2
\;\le\;
2\,\bm{d}_{\mathrm{res},k}^{\!\top}\bm{u}_{\mathrm{RL},k}.
}\end{equation}

\textbf{Step 2 --- Cauchy--Schwarz, with the known bound.} The inner
product on the right of Eq.~\ref{eq-squared} depends on the unmeasurable
\(\bm{d}_{\mathrm{res},k}\). Writing
\(\bm{d}_{\mathrm{res},k}^{\!\top}\bm{u}_{\mathrm{RL},k} =
\|\bm{d}_{\mathrm{res},k}\|\,\|\bm{u}_{\mathrm{RL},k}\|\cos\theta\) with
\(\theta \in [0,\pi]\) the angle between the two vectors, the
Cauchy--Schwarz inequality gives \(\cos\theta \le 1\), and the residual
budget Eq.~\ref{eq-budget} gives
\(\|\bm{d}_{\mathrm{res},k}\| \le \bar{d}_{\mathrm{res}}\). Together,
\begin{equation}\protect\phantomsection\label{eq-cs}{
\bm{d}_{\mathrm{res},k}^{\!\top}\bm{u}_{\mathrm{RL},k}
\;\le\;
\|\bm{d}_{\mathrm{res},k}\|\,\|\bm{u}_{\mathrm{RL},k}\|
\;\le\;
\bar{d}_{\mathrm{res}}\,\|\bm{u}_{\mathrm{RL},k}\|.
}\end{equation}

\textbf{Step 3 --- Solve for \(\|\bm{u}_{\mathrm{RL},k}\|\).}
Substituting Eq.~\ref{eq-cs} into Eq.~\ref{eq-squared}, \[
\|\bm{u}_{\mathrm{RL},k}\|^2
\;\le\;
2\,\bar{d}_{\mathrm{res}}\,\|\bm{u}_{\mathrm{RL},k}\|.
\] Dividing both sides by \(\|\bm{u}_{\mathrm{RL},k}\| > 0\) (the case
\(\bm{u}_{\mathrm{RL},k}=\bm{0}\) is trivial) yields the ceiling.

\begin{theorem}[]\protect\hypertarget{thm-ceiling}{}\label{thm-ceiling}

\textbf{Theorem 2 (ISS action-magnitude ceiling).} Suppose
Assumption~\ref{asm-lyap} holds and the baseline satisfies
Theorem~\ref{thm-baseline-iss}. For any step \(k\) at which the
preservation criterion Eq.~\ref{eq-criterion} is demanded, the RL action
magnitude must satisfy
\begin{equation}\protect\phantomsection\label{eq-ceiling}{
\|\bm{u}_{\mathrm{RL},k}\|
\;\le\;
2\,\bar{d}_{\mathrm{res}}
\;=\;
\frac{2\,c_x\,r^2}{\gamma_0}.
}\end{equation}

\end{theorem}

\subsection{The clipping law}\label{sec-cliplaw}

The ceiling Eq.~\ref{eq-ceiling} is enforced as a \emph{stateless,
radial} saturation layer interposed between the actor network and the
joint command. If the raw output already satisfies the ceiling, it
passes through unchanged; otherwise it is scaled radially back onto the
boundary of the admissible ball:

\begin{definition}[]\protect\hypertarget{def-clip}{}\label{def-clip}

\textbf{Definition 2 (ISS clipping law).} The clipped RL torque applied
to the cascade is
\begin{equation}\protect\phantomsection\label{eq-clip}{
\boxed{\quad
\bm{u}_{\mathrm{RL},k}^{\mathrm{clip}}
\;=\;
\bm{u}_{\mathrm{RL},k}\cdot
\min\!\Bigl[\,1,\;
\frac{2\,c_x\,r^2}{\gamma_0\,\|\bm{u}_{\mathrm{RL},k}\|}
\Bigr].
\quad}
}\end{equation}

\end{definition}

Two properties of Eq.~\ref{eq-clip} deserve emphasis:

\begin{enumerate}
\def\labelenumi{\arabic{enumi}.}
\item
  \textbf{Direction-preserving.} The scaling is radial, so a well-aimed
  but over-strong correction is \emph{softened}, never
  \emph{redirected}. This is important because the policy's reward
  function rewards alignment of \(\bm{u}_{\mathrm{RL}}\) with
  \(\bm{d}_{\mathrm{res}}\); per-component clamping would destroy that
  alignment.
\item
  \textbf{Stateless.} The clip depends only on the constants
  \(c_x,\gamma_0,r\) (fixed once and for all) and on the current
  \(\bm{u}_{\mathrm{RL},k}\). It does not depend on the state
  \(\bm{x}_k\) or on the residual \(\bm{d}_{\mathrm{res},k}\). This
  makes the clip cheap on hardware and trivial to implement.
\end{enumerate}

\subsection{Geometric interpretation and honest caveat}\label{sec-geom}

Eq.~\ref{eq-ceiling} is a \emph{necessary} condition for
Eq.~\ref{eq-criterion} to be satisfiable; it is \textbf{not} a blanket
\emph{sufficient} condition. Reading Eq.~\ref{eq-squared} exactly, the
preservation criterion holds if and only if
\(\|\bm{u}_{\mathrm{RL},k}\| \le 2\,\|\bm{d}_{\mathrm{res},k}\|\,\cos\theta\).
The largest value the right-hand side can take, over all admissible
residuals and all angles, is \(2\,\bar{d}_{\mathrm{res}}\), attained
only when the action is \emph{perfectly aligned} with the residual at
the worst-case magnitude. Thus \(2\,\bar{d}_{\mathrm{res}}\) is the
\textbf{outer envelope}: any action exceeding it cannot satisfy
Eq.~\ref{eq-criterion} under \emph{any} alignment, and is therefore
provably unsafe; the clip removes exactly these actions.

Inside the envelope, alignment is the policy's job, not the clip's. What
the clip guarantees, even in the adversarial worst case where the
clipped action is \emph{anti-aligned} with the residual, is a hard upper
bound on the effective residual via the triangle inequality: \[
\|\bm{d}_{\mathrm{eff},k}\|
\;\le\;
\|\bm{d}_{\mathrm{res},k}\| + \|\bm{u}_{\mathrm{RL},k}^{\mathrm{clip}}\|
\;\le\;
\bar{d}_{\mathrm{res}} + 2\,\bar{d}_{\mathrm{res}}
\;=\;
3\,\bar{d}_{\mathrm{res}}.
\] Plugging \(3\,\bar{d}_{\mathrm{res}}\) back into the ISS inequality
Eq.~\ref{eq-iss} in place of \(\|\bm{d}_{\mathrm{res}}\|\) shows that
the closed loop remains ISS with an enlarged---but still
finite---ultimate radius \(\sqrt{3}\,r\). \textbf{The cage never opens.}

\begin{tcolorbox}[enhanced jigsaw, arc=.35mm, bottomrule=.15mm, bottomtitle=1mm, breakable, colback=white, colbacktitle=quarto-callout-tip-color!10!white, colframe=quarto-callout-tip-color-frame, coltitle=black, left=2mm, leftrule=.75mm, opacityback=0, opacitybacktitle=0.6, rightrule=.15mm, title=\textcolor{quarto-callout-tip-color}{\faLightbulb}\hspace{0.5em}{What the clip guarantees, in one sentence}, titlerule=0mm, toprule=.15mm, toptitle=1mm]

If the RL policy aligns with the residual (its trained behaviour), the
error settles inside the target radius \(r\); if the policy fails
arbitrarily, the error stays inside \(\sqrt{3}\,r\).

\end{tcolorbox}

\section{Practical Estimation of the ISS Constants}\label{sec-estim}

The clipping law Eq.~\ref{eq-clip} contains three constants: \(c_x\),
\(\gamma_0\), and \(r\). This section explains how each is obtained.

\subsection{\texorpdfstring{Convergence coefficient
\(c_x\)}{Convergence coefficient c\_x}}\label{sec-estim-cx}

By Assumption (A2), \(c_x\) is the largest scalar satisfying
\(\ell(\bm{x},\kappa_N(\bm{x})) \ge c_x \|\bm{x}\|^2\). Equivalently, by
Lemma~\ref{lem-descent}, \(c_x\) is the largest scalar consistent with
the descent property
\begin{equation}\protect\phantomsection\label{eq-cx-defn}{
\frac{V_N^*(\bm{x}_k) - V_N^*(\bm{x}_{k+1}^{\mathrm{nom}})}
     {\|\bm{x}_k\|^2}
\;\ge\; c_x
\quad\text{for all nominal } \bm{x}_k\in\mathcal{X}.
}\end{equation}

A nominal trajectory is one in which the disturbance is small or absent;
on real hardware we approximate ``nominal'' by \emph{quiet} operating
conditions (no intentional disturbance, no external excitation, no
payload changes). The empirical estimator is then

\begin{definition}[]\protect\hypertarget{def-cx-emp}{}\label{def-cx-emp}

\textbf{Definition 3 (Empirical convergence coefficient).} Let
\(\{(\bm{x}_k, V_N^*(\bm{x}_k))\}_{k=1}^{K}\) be a sequence of
state--cost pairs logged from quiet closed-loop operation. The empirical
convergence coefficient is the lower 5th percentile,
\begin{equation}\protect\phantomsection\label{eq-cx-emp}{
\hat{c}_x \;=\;
\mathrm{percentile}_{5}\!\Biggl\{\;
\frac{V_N^*(\bm{x}_k) - V_N^*(\bm{x}_{k+1})}
     {\|\bm{x}_k\|^2}
\;:\;
\|\bm{x}_k\| > \epsilon,\;\;
V_N^*(\bm{x}_k) > V_N^*(\bm{x}_{k+1})
\;\Biggr\},
}\end{equation} where \(\epsilon\) is a small floor
(e.g.~\(10^{-3}\,\)m) preventing the ratio from exploding near the
origin.

\end{definition}

The use of the 5th percentile rather than the strict minimum trades off
two competing concerns: the strict minimum is contaminated by isolated
numerical outliers (e.g.~MPC solves that hit the iteration cap), while
the median over-estimates the descent and would violate the lower-bound
interpretation for a non-trivial fraction of samples. The 5th percentile
gives a \emph{robust lower envelope} that is safe in the ISS sense
(i.e.~a valid lower bound for \(\ge 95\%\) of samples).

This procedure can be applied identically in simulation and on real
hardware. We recommend taking
\(\hat{c}_x = \min(\hat{c}_x^{\mathrm{sim}},\hat{c}_x^{\mathrm{real}})\)
to absorb the discrepancy caused by unmodelled friction on the physical
plant.

\subsection{\texorpdfstring{Disturbance gain
\(\gamma_0\)}{Disturbance gain \textbackslash gamma\_0}}\label{sec-estim-gamma}

By Assumption (A3) and the construction of Eq.~\ref{eq-lip},
\(\gamma_0\) is the smallest scalar that upper-bounds the cost change
induced by a state displacement of the magnitude of
\(\bm{d}_{\mathrm{res},k}\). Re-arranging the ISS inequality
Eq.~\ref{eq-iss},
\begin{equation}\protect\phantomsection\label{eq-gamma-defn}{
\gamma_0 \;\ge\;
\frac{-\bigl[V_N^*(\bm{x}_{k+1}) - V_N^*(\bm{x}_k)\bigr]
       \;+\; c_x\,\|\bm{x}_k\|^2}
     {\|\bm{d}_{\mathrm{res},k}\|}
\quad\text{whenever } \|\bm{d}_{\mathrm{res},k}\| > 0.
}\end{equation}

The numerator is computable from logged \(V_N^*\) and \(\|\bm{x}_k\|\);
the denominator requires the \emph{clean ground-truth residual}. This is
the critical practical asymmetry:

\begin{itemize}
\tightlist
\item
  \textbf{In simulation}, \(\bm{d}_{\mathrm{res},k}\) is computable
  exactly from \(\bm{d}_{\mathrm{true},k}^{\mathrm{sim}} -
  \bm{d}_{\mathrm{filter},k}\), where the ``true'' disturbance is
  obtained by inverse dynamics with the ground-truth joint acceleration
  provided by the physics engine.
\item
  \textbf{On real hardware}, the joint acceleration is reconstructed by
  finite differences, which injects numerical noise into the residual
  estimate. The empirical fit Eq.~\ref{eq-gamma-defn} would
  \emph{under-estimate} \(\gamma_0\) (because the noisy denominator is
  inflated), which is the unsafe direction.
\end{itemize}

The remedy is to estimate \(\gamma_0\) \emph{in simulation only}, then
transfer the value to real hardware. The transfer is exact because
\(\gamma_0\) is a property of the cost function \(V_N^*\), and \(V_N^*\)
is the \emph{same mathematical function} in sim and real (same NMPC
weights, same URDF, same horizon, same sample time).

\begin{definition}[]\protect\hypertarget{def-gamma-emp}{}\label{def-gamma-emp}

\textbf{Definition 4 (Empirical disturbance gain).} Let
\(\{(\bm{x}_k, V_N^*(\bm{x}_k), \bm{d}_{\mathrm{res},k}^{\mathrm{sim}})\}_{k=1}^{K}\)
be logged from a baseline simulation run with the full disturbance
curriculum active. The empirical disturbance gain is the 95th
percentile, \[
\hat{\gamma}_0 \;=\;
\mathrm{percentile}_{95}\!\Biggl\{\;
\frac{-\Delta V_k + \hat{c}_x\,\|\bm{x}_k\|^2}
     {\|\bm{d}_{\mathrm{res},k}^{\mathrm{sim}}\|}
\;:\;
\|\bm{d}_{\mathrm{res},k}^{\mathrm{sim}}\| > \delta,
\;\;\text{(ratio)} > 0
\;\Biggr\},
\] where \(\Delta V_k = V_N^*(\bm{x}_{k+1}) - V_N^*(\bm{x}_k)\) and
\(\delta\) is a small floor preventing division by near-zero residuals.

\end{definition}

The 95th percentile---rather than the strict maximum---gives a
\emph{robust upper envelope} in the same way the 5th percentile gave a
robust lower envelope for \(\hat{c}_x\).

\subsection{\texorpdfstring{Target radius
\(r\)}{Target radius r}}\label{sec-estim-r}

The radius \(r\) is the \emph{only} design quantity in the clipping law;
it is not measured but \emph{selected}. It encodes the practical
tracking accuracy the user is willing to certify. Three observations
guide the choice:

\begin{enumerate}
\def\labelenumi{\arabic{enumi}.}
\item
  \textbf{The clip ceiling scales as \(r^2\).} Halving \(r\) shrinks the
  ceiling by a factor of four; doubling \(r\) enlarges it sixteen-fold.
  Aggressive (small) \(r\) yields a tighter guarantee but allows the RL
  layer almost no contribution; permissive (large) \(r\) enlarges the
  guarantee region but lets the policy be more active.
\item
  \textbf{\(r\) should be consistent with what the baseline can already
  achieve.} If the baseline NMPC--DOB tracking error is observed to
  settle around \(\|\bm{x}_k\| \approx r_{\mathrm{base}}\) in steady
  state, then choosing \(r \ll r_{\mathrm{base}}\) certifies an
  unattainable guarantee and forces the clip to be uselessly tight. A
  pragmatic choice is \(r \approx r_{\mathrm{base}}\), or slightly
  smaller if the RL layer is expected to close the gap.
\item
  \textbf{\(r\) does not appear in the ISS certificate of
  Theorem~\ref{thm-baseline-iss} itself.} It only appears in the
  \emph{budget} Eq.~\ref{eq-budget} and the \emph{ceiling}
  Eq.~\ref{eq-ceiling}. ISS holds independently of \(r\); the choice of
  \(r\) only specifies \emph{which} radius the budget targets.
\end{enumerate}

\subsection{Hybrid sim--real workflow}\label{sec-estim-flow}

Putting the three estimators together:

{\def\LTcaptype{none} % do not increment counter
\begin{longtable}[]{@{}
  >{\raggedright\arraybackslash}p{(\linewidth - 4\tabcolsep) * \real{0.2211}}
  >{\raggedright\arraybackslash}p{(\linewidth - 4\tabcolsep) * \real{0.3263}}
  >{\raggedright\arraybackslash}p{(\linewidth - 4\tabcolsep) * \real{0.4526}}@{}}
\toprule\noalign{}
\begin{minipage}[b]{\linewidth}\raggedright
Constant
\end{minipage} & \begin{minipage}[b]{\linewidth}\raggedright
Where measured
\end{minipage} & \begin{minipage}[b]{\linewidth}\raggedright
Why
\end{minipage} \\
\midrule\noalign{}
\endhead
\bottomrule\noalign{}
\endlastfoot
\(c_x\) & both, take \(\min\) & property of NMPC + plant; real may be
smaller due to friction \\
\(\gamma_0\) & simulation only & property of \(V_N^*\) alone; transfers
exactly to real \\
\(r\) & design choice & not measured \\
\end{longtable}
}

The final clip ceiling used at run-time is
\begin{equation}\protect\phantomsection\label{eq-final-clip}{
\|\bm{u}_{\mathrm{RL}}\|
\;\le\;
\frac{2\,\hat{c}_x\,r^2}{\hat{\gamma}_0},
\qquad
\hat{c}_x = \min(\hat{c}_x^{\mathrm{sim}}, \hat{c}_x^{\mathrm{real}}),
\quad
\hat{\gamma}_0 = \hat{\gamma}_0^{\mathrm{sim}}.
}\end{equation}

Because the strict minimum of \(c_x\) and the strict maximum form of
\(\gamma_0\) are \emph{both} taken in the conservative direction, all
estimation error is absorbed by a small reduction in the clip ceiling.
Stability is therefore preserved under estimation error; only the RL
layer's freedom is reduced. This is the desired robustness property of
the proposed certificate.

\section{Summary}\label{sec-summary}

Under the standing assumptions Assumption~\ref{asm-lyap}, the NMPC--DOB
cascade is ISS with respect to the residual disturbance
(Theorem~\ref{thm-baseline-iss}), with quadratic dissipation rate
\(c_x\) and linear disturbance gain \(\gamma_0\). For any target radius
\(r > 0\), the admissible residual budget
\(\bar{d}_{\mathrm{res}} = c_x r^2/\gamma_0\) specifies the largest
residual the baseline can absorb while keeping the tracking error inside
\(r\). The residual RL layer is admissible at step \(k\) if it does not
enlarge the residual it is meant to cancel; this preservation criterion
reduces, via the Cauchy--Schwarz inequality, to the magnitude ceiling
\(\|\bm{u}_{\mathrm{RL},k}\| \le 2 c_x r^2/\gamma_0\)
(Theorem~\ref{thm-ceiling}) and the radial clipping law
Eq.~\ref{eq-clip}. Provably-unsafe actions---those whose magnitude
exceeds \(2\bar{d}_{\mathrm{res}}\) and cannot satisfy the preservation
criterion under any alignment---are removed by construction. Inside the
envelope, even adversarial misalignment of the clipped action leaves the
closed loop ISS with enlarged ultimate radius \(\sqrt{3}\,r\); the cage
never opens. The three constants \(c_x,
\gamma_0, r\) are obtained via Definition~\ref{def-cx-emp},
Definition~\ref{def-gamma-emp}, and the design rule of
Section~\ref{sec-estim-r}; the resulting hybrid sim--real workflow
Eq.~\ref{eq-final-clip} is robust to estimation error in the
conservative direction.

\section{Symbol Reference}\label{symbol-reference}

{\def\LTcaptype{none} % do not increment counter
\begin{longtable}[]{@{}
  >{\raggedright\arraybackslash}p{(\linewidth - 2\tabcolsep) * \real{0.3824}}
  >{\raggedright\arraybackslash}p{(\linewidth - 2\tabcolsep) * \real{0.6176}}@{}}
\toprule\noalign{}
\begin{minipage}[b]{\linewidth}\raggedright
Symbol
\end{minipage} & \begin{minipage}[b]{\linewidth}\raggedright
Meaning
\end{minipage} \\
\midrule\noalign{}
\endhead
\bottomrule\noalign{}
\endlastfoot
\(\bm{x}_k \in \mathbb{R}^{12}\) & Tracking-error state (6 position
errors + 6 velocity errors) \\
\(\bm{u}_{\mathrm{MPC},k}\) & Nominal NMPC torque,
\(\kappa_N(\bm{x}_k)\) \\
\(\bm{d}_{\mathrm{filter},k}\) & DOB low-pass-filtered disturbance
estimate \\
\(\bm{u}_{\mathrm{RL},k}\) & Residual RL torque (pre-clip) \\
\(\bm{u}_{\mathrm{RL},k}^{\mathrm{clip}}\) & RL torque after ISS
clipping Eq.~\ref{eq-clip} \\
\(\bm{d}_{\mathrm{res},k}\) & Residual disturbance (DOB out-of-band
leftover) \\
\(\bm{d}_{\mathrm{eff},k}\) & Effective residual under RL,
\(\bm{d}_{\mathrm{res}}-\bm{u}_{\mathrm{RL}}\) \\
\(V_N^*(\bm{x})\) & NMPC optimal cost (Lyapunov candidate) \\
\(\kappa_N(\bm{x})\) & NMPC feedback law (first element of optimal
sequence) \\
\(\ell(\bm{x},\bm{u})\) & NMPC per-step stage cost \\
\(c_x > 0\) & Convergence coefficient, stage-cost lower bound (A2) \\
\(\gamma_0 > 0\) & Disturbance gain, local Lipschitz constant of
\(V_N^*\) (A3) \\
\(r > 0\) & Target practical convergence radius \\
\(\bar{d}_{\mathrm{res}} = c_x r^2/\gamma_0\) & Admissible residual
budget (Definition~\ref{def-budget}) \\
\(2\bar{d}_{\mathrm{res}} = 2c_x r^2/\gamma_0\) & ISS action ceiling
(Theorem~\ref{thm-ceiling}) \\
\end{longtable}
}

\section*{References}\label{references}
\addcontentsline{toc}{section}{References}

\protect\phantomsection\label{refs}

\begin{tcolorbox}[enhanced jigsaw, arc=.35mm, bottomrule=.15mm, bottomtitle=1mm, breakable, colback=white, colbacktitle=quarto-callout-note-color!10!white, colframe=quarto-callout-note-color-frame, coltitle=black, left=2mm, leftrule=.75mm, opacityback=0, opacitybacktitle=0.6, rightrule=.15mm, title=\textcolor{quarto-callout-note-color}{\faInfo}\hspace{0.5em}{Suggested bibliography entries (\texttt{refs.bib})}, titlerule=0mm, toprule=.15mm, toptitle=1mm]

\begin{Shaded}
\begin{Highlighting}[]
\VariableTok{@article}\NormalTok{\{}\OtherTok{mayne2000}\NormalTok{,}
  \DataTypeTok{author}\NormalTok{  = \{D. Q. Mayne and J. B. Rawlings and C. V. Rao and P. O. M. Scokaert\},}
  \DataTypeTok{title}\NormalTok{   = \{Constrained model predictive control: stability and optimality\},}
  \DataTypeTok{journal}\NormalTok{ = \{Automatica\},}
  \DataTypeTok{volume}\NormalTok{  = \{36\},}
  \DataTypeTok{number}\NormalTok{  = \{6\},}
  \DataTypeTok{pages}\NormalTok{   = \{789{-}{-}814\},}
  \DataTypeTok{year}\NormalTok{    = \{2000\}}
\NormalTok{\}}

\VariableTok{@book}\NormalTok{\{}\OtherTok{rawlings2017}\NormalTok{,}
  \DataTypeTok{author}\NormalTok{    = \{J. B. Rawlings and D. Q. Mayne and M. M. Diehl\},}
  \DataTypeTok{title}\NormalTok{     = \{Model Predictive Control: Theory, Computation, and Design\},}
  \DataTypeTok{edition}\NormalTok{   = \{2nd\},}
  \DataTypeTok{publisher}\NormalTok{ = \{Nob Hill Publishing\},}
  \DataTypeTok{year}\NormalTok{      = \{2017\}}
\NormalTok{\}}

\VariableTok{@article}\NormalTok{\{}\OtherTok{jiang2001}\NormalTok{,}
  \DataTypeTok{author}\NormalTok{  = \{Z.{-}P. Jiang and Y. Wang\},}
  \DataTypeTok{title}\NormalTok{   = \{Input{-}to{-}state stability for discrete{-}time nonlinear systems\},}
  \DataTypeTok{journal}\NormalTok{ = \{Automatica\},}
  \DataTypeTok{volume}\NormalTok{  = \{37\},}
  \DataTypeTok{number}\NormalTok{  = \{6\},}
  \DataTypeTok{pages}\NormalTok{   = \{857{-}{-}869\},}
  \DataTypeTok{year}\NormalTok{    = \{2001\}}
\NormalTok{\}}

\VariableTok{@book}\NormalTok{\{}\OtherTok{khalil2002}\NormalTok{,}
  \DataTypeTok{author}\NormalTok{    = \{H. K. Khalil\},}
  \DataTypeTok{title}\NormalTok{     = \{Nonlinear Systems\},}
  \DataTypeTok{edition}\NormalTok{   = \{3rd\},}
  \DataTypeTok{publisher}\NormalTok{ = \{Prentice Hall\},}
  \DataTypeTok{year}\NormalTok{      = \{2002\}}
\NormalTok{\}}

\VariableTok{@article}\NormalTok{\{}\OtherTok{sontag1995}\NormalTok{,}
  \DataTypeTok{author}\NormalTok{  = \{E. D. Sontag and Y. Wang\},}
  \DataTypeTok{title}\NormalTok{   = \{On characterizations of the input{-}to{-}state stability property\},}
  \DataTypeTok{journal}\NormalTok{ = \{Systems }\CharTok{\textbackslash{}\&}\NormalTok{ Control Letters\},}
  \DataTypeTok{volume}\NormalTok{  = \{24\},}
  \DataTypeTok{pages}\NormalTok{   = \{351{-}{-}359\},}
  \DataTypeTok{year}\NormalTok{    = \{1995\}}
\NormalTok{\}}
\end{Highlighting}
\end{Shaded}

\end{tcolorbox}




\end{document}
